Die Erstellung eines qualitativ hochwertigen, repräsentativen und ausreichend großen parallelen Korpus bildet das Fundament für jedes datengetriebene Übersetzungssystem. Wie im vorangegangenen Kapitel dargelegt, leidet die maschinelle Textvereinfachung (ATS) für die deutsche Leichte Sprache unter einer eklatanten Datenknappheit. Während bestehende Ressourcen wie das Korpus von \citet{battisti2020corpus} einen wertvollen Grundstein legten, weisen sie für das Training moderner Seq2Seq- und Präferenzoptimierungsmodelle (wie DPO) entscheidende Limitierungen auf: Entweder ist ihr paralleler Anteil mit wenigen hundert Dokumenten zu gering, oder sie fokussieren sich, wie das DEplain-Korpus von \citet{stodden2023deplain}, primär auf Einfache Sprache (Plain German), was sich linguistisch und regulatorisch signifikant von der Leichten Sprache unterscheidet. Einen wichtigen methodischen Impuls zur automatisierten Datengewinnung lieferten \citet{toborek2023new} mit dem \textit{New Aligned Simple German Corpus}. Ihr Ansatz demonstrierte, wie parallele Webinhalte über seitenspezifische Selektoren und Sprachumschalter im Internet automatisiert identifiziert und gescraped werden können. Diese Vorgehensweise bildete den primären Ausgangspunkt für die in dieser Arbeit entwickelte Crawling- und Extraktions-Pipeline. Allerdings fokussieren Toborek et al. primär auf ein inhärent fehleranfälliges Satz-Alignment und differenzieren nicht strikt zwischen Leichter und Einfacher Sprache.

Um diese Forschungslücke zu schließen, wurde im Rahmen dieser Arbeit eine eigene, automatisierte Pipeline zur Akquise und Aufbereitung von Trainingsdaten entwickelt, welche die Scraping-Prinzipien von \citet{toborek2023new} als methodisches Fundament aufgreift, auf 11 heterogene Domänen ausweitet und zu einer artikelbasierten, qualitätsgesicherten Korpus-Pipeline weiterentwickelt. Da diese Vorarbeit mittlerweile rund drei Jahre zurückliegt und im deutschsprachigen Web in diesem Zeitraum zahlreiche neue Informationsangebote sowie redaktionelle Portale in Leichter Sprache entstanden sind, wurde die Entscheidung getroffen, einen vollständig neuen, zeitgemäßen und breiter aufgestellten Korpus aufzubauen.

Ein zentrales methodisches Charakteristikum der in dieser Arbeit entwickelten Pipeline und der nachfolgenden Modellierung ist die bewusste Entscheidung für ein \textbf{artikelbasiertes} (bzw. dokumentenweites) Training sowohl der Komplexitätsmetrik als auch des Übersetzungsmodells. In der bisherigen Forschung wurden Vereinfachungssysteme primär auf Satzebene trainiert und evaluiert, was jedoch erhebliche Defizite aufweist:
\begin{enumerate}
    \item \textbf{Kontext- und Kohärenzverlust:} Satzbasierte Ansätze isolieren Sätze aus ihrem diskursiven Zusammenhang. Dadurch gehen anaphorische Bezüge, der rote Faden sowie die globale Kohärenz verloren. Zudem können makrostrukturelle Vereinfachungen wie die Umstrukturierung ganzer Absätze oder das Einfügen von erklärenden Hintergrundinformationen – welche im Konzept der Leichten Sprache essenziell sind – auf Satzebene nicht gelernt werden.
    \item \textbf{Verrauschte Satz-Alignments:} Die Modellierung auf Satzebene setzt voraus, dass präzise Alignments vorliegen. Wie \citet{stodden2023deplain} und \citet{stajner2018cats} zeigen, führt die asymmetrische Natur der Vereinfachung (z.~B. n:m-Alignments und Satzteilungen) dazu, dass automatische Satz-Aligner entweder stark fehlerhafte Paare erzeugen oder einen massiven Informationsverlust durch das Verwerfen nicht-linearer Übergänge in Kauf nehmen müssen.
    \item \textbf{Verlust struktureller Metadaten:} Wie \citet{battisti2020corpus} nachwiesen, spielen visuelle Layoutmerkmale, Absatzstrukturen und die Gliederung auf Dokumentenebene eine entscheidende Rolle für das Verständnis von Leichter Sprache. Ein satzbasiertes Alignment eliminiert diese Strukturen vollständig.
\end{enumerate}
Indem wir den gesamten Artikel als Basiseinheit wählen, kann das Übersetzungsmodell lernen, globale Kürzungen vorzunehmen und Texte kohärent umzuformulieren. Gleichzeitig erlaubt dies der Komplexitätsmetrik, die Verständlichkeit eines Textes im Kontext des gesamten Dokuments statt isolierter Sätze zu bewerten. Die methodischen Entscheidungen dieser Pipeline spiegeln dabei direkt die in Abschnitt~\ref{sec:maschinelle_textvereinfachung_alignment} und \ref{sec:evaluierungsmetriken} diskutierten Herausforderungen des Forschungsstandes wider:
\begin{itemize}
    \item \textbf{Multi-Quellen-Crawling \& Datenakquise (Abschnitt~3.1):} Aufbauend auf den Scraping- und Alignment-Heuristiken des \textit{New Aligned Simple German Corpus} \citep{toborek2023new} als Ausgangspunkt wurde ein modularer Scraper für eine Vielzahl deutschsprachiger Webportale implementiert. Dies stellt sicher, dass sowohl behördliche Texte (im Kontext der BITV 2.0) als auch journalistische Beiträge und Alltagstexte über 11 heterogene Quellen hinweg automatisiert erfasst werden.
    \item \textbf{Alignment \& Das 1:n-Problem (Abschnitt~3.2):} Wie in Abschnitt~\ref{sec:maschinelle_textvereinfachung_alignment} erläutert, übersetzen menschliche Übersetzer Leichte Sprache selten in Form von 1:1-Satzentsprechungen. Häufiger sind Satzaufspaltungen (Sentence Splitting) und starke Komprimierungen. Daher wird eine dynamische Alignment-Strategie benötigt, die über klassische Satzähnlichkeiten hinausgeht und semantische Einbettungen nutzt, um asymmetrische Satzbeziehungen (1:n) zuverlässig abzubilden.
    \item \textbf{Quality Assurance \& Similarity Sweet-Spot (Abschnitt~3.3):} Um das von \citet{stajner2018cats} beschriebene Rauschen in den ausgerichteten Daten zu minimieren, wird eine filterbasierte Qualitätssicherung eingeführt. Diese zielt darauf ab, den "`Similarity Sweet-Spot"' zu identifizieren: Einerseits müssen zu unähnliche Satzpaare (die auf fehlerhaften Alignments beruhen) ausgeschlossen werden; andererseits müssen auch identische oder nahezu identische Satzpaare (die keinen Vereinfachungseffekt aufweisen) herausgefiltert werden, um dem Modell das Erlernen trivialer Kopieroperationen zu erschweren.
    \item \textbf{Datenbereinigung \& Normalisierung (Abschnitt~3.4):} Da Texte aus unterschiedlichen Quellen stark variierende HTML-Strukturen, typografische Besonderheiten (z.~B. die Verwendung von Bindestrichen oder Mediopunkten zur Worttrennung) und Layout-Elemente aufweisen, ist eine sorgfältige Vorverarbeitung essenziell.
    \item \textbf{Korpus-Statistiken \& Lexikalische Diversität (Abschnitt~3.5):} Abschließend wird das resultierende Korpus quantitativ und qualitativ evaluiert. Hierbei wird analysiert, inwiefern die bereinigten und ausgerichteten Daten die theoretischen Anforderungen der Leichten Sprache erfüllen und welche lexikalische Vielfalt sie aufweisen.
\end{itemize}

Dieses Kapitel beschreibt den detaillierten Ablauf dieser Schritte und dokumentiert die Designentscheidungen bei der Konstruktion des finalen Datensatzes.

\subsection{Multi-Quellen-Crawling \& Datenakquise}
\label{sec:crawling_datenakquise}

Die automatisierte Akquise paralleler Textdaten für die deutsche Leichte Sprache stellt aufgrund der heterogenen Weblandschaft und des Fehlens zentraler Repositorien eine erhebliche methodische Herausforderung dar. Während im englischsprachigen Raum Plattformen wie die \textit{Simple English Wikipedia} als quasi-standardisierte Datenquelle dienen, existiert im deutschsprachigen Raum kein vergleichbares, monolithisches Korpus. Parallele Sprachversionen verteilen sich stattdessen über behördliche Portale, öffentlich-rechtliche Rundfunkanstalten, Verbände der Behindertenhilfe und journalistische Magazine. 

Um diesen heterogenen Datenbestand automatisiert und reproduzierbar zu erschließen, greift die in dieser Arbeit entwickelte Pipeline die grundlegenden Scraping-Heuristiken des \textit{New Aligned Simple German Corpus} von \citet{toborek2023new} als methodischen Ausgangspunkt auf. Während Toborek et al. ihr Verfahren jedoch primär auf die Extraktion unzusammenhängender Satzfragmente ausrichteten, wird der Ansatz hier zu einer modularen, zweistufigen Dokumenten-Crawling-Architektur erweitert. Diese trennt die logische URL-Identifikation strikt von der inhaltlichen DOM-Extraktion:

\begin{enumerate}
    \item \textbf{Stufe 1: URL-Discovery \& Alignment-Heuristiken (\texttt{scripts/data\_collection/crawl\_scraper/}):} In diesem Schritt navigieren spezialisierte Crawler durch die Zielportale, identifizieren Artikel in Leichter Sprache (LS) und ermitteln über seitenspezifische Verknüpfungsmechanismen die korrespondierende URL der Alltagssprache (AS). Das Ergebnis jeder Quelle wird als JSON-Struktur standardisierter URL-Paare abgelegt.
    \item \textbf{Stufe 2: Selektive DOM-Content-Extraktion (\texttt{scripts/data\_collection/corpus\_scrapers/}):} Aufbauend auf den identifizierten URL-Paaren lädt die zweite Stufe die HTML-Dokumente asynchron herunter. Anstatt den gesamten Seiteninhalt ungefiltert zu erfassen, extrahiert ein zielgerichteter DOM-Parser ausschließlich den redaktionellen Fließtext über semantische HTML5-Container (\texttt{<main>}, \texttt{<article>}, \texttt{div.paragraph}). Navigationsleisten, Kopf- und Fußzeilen, Cookie-Banner sowie interaktive Schaltflächen werden dadurch bereits auf Extraktionsebene systematisch eliminiert.
\end{enumerate}

\subsubsection{Domänen- und Quellenabdeckung}
Um eine breite stilistische und thematische Varianz zu gewährleisten und eine Überanpassung an ein einzelnes redaktionelles Regelwerk zu verhindern, wurden 12 frei zugängliche deutschsprachige Webquellen aus drei gesellschaftlich relevanten Domänen erschlossen:

\begin{itemize}
    \item \textbf{Öffentliche Verwaltung \& Kommunalportale:} Dieser Sektor besitzt besondere rechtliche Relevanz, da öffentliche Stellen gemäß der Barrierefreie-Informationstechnik-Verordnung (BITV 2.0) gesetzlich verpflichtet sind, wesentliche Informationen in Leichter Sprache bereitzustellen \citep{bund2026barrierefreiheit}. Erfasst wurden die Stadtportale \textit{Hannover.de} (welches aufgrund seines Umfangs als Referenzkorpus dient), \textit{Hamburg.de}, \textit{Koeln.de}, \textit{Stuttgart.de}, \textit{Wiesbaden.de} sowie das Portal des \textit{Main-Taunus-Kreises}.
    \item \textbf{Medien \& Journalismus:} Zur Abdeckung tagesaktueller Berichterstattung und nachrichtlicher Diskursformen umfasst das Korpus die Angebote des Mitteldeutschen Rundfunks (\textit{MDR Nachrichten in Leichter Sprache}), der tageszeitung (\textit{taz}) sowie des Wirtschaftsmagazins \textit{brand eins}.
    \item \textbf{Gesundheit, Soziales \& Politik:} Zur Erfassung medizinischer und sozialrechtlicher Fachtexte wurden das Portal der \textit{Apotheken Umschau} (Einfache Sprache), die Veröffentlichungen des \textit{Beauftragten der Bundesregierung für die Belange von Menschen mit Behinderungen} sowie das Bildungsportal \textit{Sozialpolitik.com} integriert.
\end{itemize}

Ergänzend zu den Webquellen wurde ein unabhängiger \textbf{Out-of-Domain-Testdatensatz (OOD)} auf Basis realer Dokumente der \textit{Lebenshilfe Schleswig-Holstein e.\,V.} aufgebaut. Dieser Datensatz umfasst behördliche Bescheide, Hausordnungen, Vereinbarungen und Prüfberichte, die von professionellen Übersetzungsbüros übertragen und von zertifizierten Prüfgruppen aus Menschen mit Lernschwierigkeiten validiert wurden. Die Extraktion erfolgte über ein spezialisiertes Parsing-Skript (\texttt{scripts/preprocessing/create\_lebenshilfe\_dataset.py}), welches proprietäre Dateiformate (\texttt{.docx}, \texttt{.rtf}, \texttt{.odt}) verarbeitet und für die spätere modellunabhängige Generalisierungsevaluation bereitstellt.

\subsubsection{Technische Extraktions- \& Verknüpfungsmechanismen}
Da Webseiten keine standardisierten Metadaten für sprachliche Vereinfachungen bereitstellen, wurden für die 12 Webquellen fünf distinkte Verknüpfungsmechanismen implementiert:

\begin{enumerate}
    \item \textbf{Canonical Link Header Exploration:} Die Stadt Hannover implementiert eine technisch vorbildliche Verknüpfung. Jeder LS-Artikel besitzt im HTML-Kopf ein Canonical-Tag (\texttt{<link rel="canonical" href="...">}), welches direkt auf den ursprünglichen Volltext in Alltagssprache verweist. Der Scraper durchsucht ausgehend von der Übersichtsseite rekursiv alle internen Pfade und bildet über diese Referenz deterministische Paare.
    \item \textbf{DOM Language Switcher Detection:} Bei Quellen wie dem \textit{Behindertenbeauftragten} (\texttt{.c-language-switch\_\_l--as}), \textit{Hamburg.de} (\texttt{.km1-icon--original-language}) oder \textit{Sozialpolitik.com} (\texttt{.underline.easy} mit \texttt{hreflang="de-DE"}) existieren dedizierte Umschalt-Schaltflächen im DOM. Der Crawler parst die LS-Seite, liest das Link-Ziel des Sprachumschalters aus und verknüpft die Dokumente.
    \item \textbf{URL-Parameter- \& Routing-Heuristiken:} Einige kommunale Content-Management-Systeme routen Sprachvarianten über URL-Query-Parameter. Bei \textit{Stuttgart.de} wird die Zielsprache über den Parameter \texttt{?sp:out=easy} gesteuert; bei \textit{Wiesbaden.de} erfolgt die Auslieferung über \texttt{?sp:easylanguage=1}. Das Alignment erfolgt hierbei durch gezieltes Setzen bzw. Entfernen des Parameters bei identischem URL-Stamm.
    \item \textbf{Kontextuelle Backlinks \& Text-Trigger:} Redaktionelle Angebote betten Links zur Ausgangssprache häufig direkt in den Fließtext oder Teaser-Boxen ein. Beim \textit{MDR} sucht der Scraper nach dem standardisierten Block \enquote{Hier können Sie diese Nachricht auch in schwerer Sprache lesen:}; die \textit{taz} verlinkt den Originalartikel im Schlusssatz innerhalb kursiver Elemente (\texttt{<em>}), während die \textit{Apotheken Umschau} Links mit dem Titelattribut \enquote{mehr Informationen} außerhalb des vereinfachten Verzeichnispfads auswertet.
    \item \textbf{Intra-Page CSS- \& Layout-Dissektion:} Das Magazin \textit{brand eins} publiziert beide Sprachfassungen gemeinsam auf einer einzigen Webseite innerhalb desselben Textcontainers (\texttt{section.textblock}). Das Alignment erfordert hier eine Stil-Dissektion: Absätze in Leichter Sprache sind im Inline-CSS rot formatiert (\texttt{color: \#fa4600} bzw. \texttt{\#ff4948}), während der schwarze Standardtext der Alltagssprache entspricht.
\end{enumerate}

Für historische oder umstrukturierte Webseiten (\textit{Stadt Köln}, \textit{Lebenshilfe Main-Taunus}, \textit{brand eins}) wurde das Scraping über die \textit{Wayback Machine} des Internet Archive realisiert. Um Verbindungsabbrüche durch serverseitige Rate-Limits zu verhindern, implementieren die Scraper eine robuste Retry-Logik mit exponentiellem Backoff (Intervall $2\,\text{s} \to 4\,\text{s} \to 8\,\text{s}$).

\subsubsection{Ergebnis der Datenakquise (Roh-Korpus)}
Durch den Einsatz dieser zielgerichteten Scraper-Infrastruktur konnten in der ersten Pipelinestufe insgesamt \textbf{1.533 parallele Artikelpaare} über alle 12 Webquellen hinweg akquiriert werden. Tabelle~\ref{tab:rohdaten_uebersicht} fasst den Rohdatenbestand nach Quellen zusammen.

\begin{table}[htbp]
\centering
\small
\caption{Übersicht des akquirierten Rohdatenbestands (Stufe 1 der Pipeline) nach Quellen.}
\label{tab:rohdaten_uebersicht}
\begin{tabularx}{\textwidth}{lrrrr}
\toprule
\textbf{Quelle} & \textbf{Artikelpaare} & \textbf{AS-Tokens} & \textbf{LS-Tokens} & \textbf{Token-Ratio (LS/AS)} \\
\midrule
Hannover.de & 808 & 871.042 & 872.115 & 1,00 \\
MDR Nachrichten & 235 & 173.210 & 98.432 & 0,57 \\
Apotheken Umschau & 161 & 451.290 & 241.108 & 0,53 \\
Behindertenbeauftragter & 60 & 51.820 & 40.612 & 0,78 \\
Hamburg.de & 57 & 67.614 & 68.012 & 1,01 \\
Stuttgart.de & 42 & 124.918 & 64.310 & 0,51 \\
Wiesbaden.de & 41 & 30.012 & 21.240 & 0,71 \\
Koeln.de & 39 & 60.320 & 92.015 & 1,53 \\
Main-Taunus-Kreis & 36 & 17.514 & 18.910 & 1,08 \\
brand eins & 32 & 22.140 & 21.218 & 0,96 \\
Sozialpolitik.com & 15 & 14.415 & 6.390 & 0,44 \\
taz (die tageszeitung) & 7 & 4.699 & 3.993 & 0,85 \\
\midrule
\textbf{Gesamt (Rohdaten)} & \textbf{1.533} & \textbf{1.868.994} & \textbf{1.468.345} & \textbf{0,79} \\
\bottomrule
\end{tabularx}
\end{table}

Insgesamt umfasst der Rohdatenbestand rund \textbf{1,87 Millionen Tokens} in der Ausgangssprache und rund \textbf{1,47 Millionen Tokens} in Leichter Sprache (gemessen mit dem Tokenisierer \texttt{cl100k\_base}). Wie die Token-Ratios in Tabelle~\ref{tab:rohdaten_uebersicht} verdeutlichen, unterscheidet sich das Längenverhalten zwischen den Quellen erheblich: Während journalistische und medizinische Quellen (MDR, Apotheken Umschau, Sozialpolitik) stark komprimieren (Ratio $\approx 0{,}44$ bis $0{,}57$), weisen kommunale Quellen (Hannover, Hamburg, Köln) ausgeglichene oder expandierende Längenverhältnisse auf (Ratio $> 1{,}00$, bei Köln bis zu $1{,}53$). 

Diese deutliche Diskrepanz in der Token-Ratio spiegelt wider, wie unterschiedlich das Regelwerk und das Zielbild von Leichter Sprache in der Praxis interpretiert und redaktionell umgesetzt werden. Während im nachrichtlichen und journalistischen Kontext die radikale inhaltliche Reduktion sowie Verdichtung auf Kernbotschaften im Vordergrund stehen, verfolgen insbesondere behördliche und kommunale Portale einen explizierenden Ansatz. Hierbei wird versucht, die vollständige administrative Information beizubehalten, wodurch kurze Fachbegriffe, formale Abläufe und juristische Sachverhalte durch detaillierte, mehrsätzige Erläuterungen und Kontextbeschreibungen aufgeschlüsselt werden. Die Token-Ratio verdeutlicht somit nicht nur statistische Längenunterschiede, sondern illustriert die Bandbreite zwischen inhaltsfilternder Zusammenfassung und vollumfänglicher syntaktisch-semantischer Erläuterung im Spektrum der Leichten Sprache.

Gleichzeitig weist der Rohdatenbestand noch erhebliche Qualitätsunterschiede auf, da fehlerhafte Verlinkungen, leere Teaser-Seiten sowie verbleibendes Web-Boilerplate enthalten sind. Die nachfolgenden Abschnitte beschreiben daher die methodischen Schritte zur Qualitätssicherung und Bereinigung.

\subsection{Alignment \& Das 1:n-Problem}
\label{sec:alignment_1_to_n}
\begin{itemize}
    \item \textbf{Ebenen des Alignments:} Gegenüberstellung von Satz- vs. Absatz- vs. Dokument-Alignment.
    \item \textbf{Asymmetrische Abbildungen (1:n und n:m):} Analyse von Satzaufteilungen (\textit{Sentence Splitting}) und syntaktischer Restrukturierung in Leichter Sprache.
    \item \textbf{Wahrung des Diskurskontextes:} Vermeidung von Kontextverlust und anaphorischen Brüchen durch dokumentenweites Pairing.
\end{itemize}

\subsection{Quality Assurance \& Similarity Sweet-Spot}
\label{sec:quality_assurance_sweet_spot}
\begin{itemize}
    \item \textbf{Long-Context Embeddings:} Einsatz von \textit{Jina Embeddings v2} (8192 Tokens Kontextfenster) zur lückenlosen Vektorisierung langer Artikel ohne Truncation-Artefakte klassischer SBERT-Modelle.
    \item \textbf{Der Similarity Sweet-Spot (Filtration):} Definition des optimalen Ähnlichkeitskorridors ($0{,}60 \le \text{sim}_{\cos} \le 0{,}98$).
    \begin{itemize}
        \item Ausschluss von Fehl-Alignments und irrelevanten Metadaten-Seiten ($\text{sim} < 0{,}60$).
        \item Ausschluss identischer Kopien und unvollständiger Übersetzungen ohne Vereinfachungseffekt ($\text{sim} > 0{,}98$).
    \end{itemize}
\end{itemize}

\subsection{Datenbereinigung \& Normalisierung}
\label{sec:datenbereinigung_normalisierung}
\begin{itemize}
    \item \textbf{Web-Boilerplate-Bereinigung:} Chirurgische Entfernung von Navigationsfragmenten, Teasern, Bildunterschriften und Radio-/Audio-Metadaten.
    \item \textbf{Typografische Normalisierung:} Behandlung von Mediopunkten ($\cdot$), Bindestrich-Komposita und Unicode-Sonderzeichen.
    \item \textbf{Glossar-Anreicherung (\textit{Enrichment}):} Extraktion und strukturierte Integration von Begriffserläuterungen aus dem \textit{Hurraki}-Wörterbuch für Leichte Sprache.
\end{itemize}

\subsection{Korpus-Statistiken \& Lexikalische Diversität}
\label{sec:korpus_statistiken_diversitaet}
\begin{itemize}
    \item \textbf{Quantitative Korpusmaße:} Artikelpaare, Token- und Satzanzahlen sowie Längenverhältnisse nach Quellen.
    \item \textbf{Lexikalische Diversität:} Vergleichende Analyse mittels Type-Token-Ratio (TTR) und Moving-Average Type-Token-Ratio (MATTR) zur Kontrolle des Längeneinflusses.
    \item \textbf{Traditionelle Lesbarkeitsindizes:} Auswertung nach Wiener Sachtextformel (WSTF 1--4), Flesch-Reading-Ease (FRE) und LIX.
    \item \textbf{Faktenerhalt \& Informationsverlust:} Bidirektionaler Named Entity Recall (NER mit spaCy \texttt{de\_core\_news\_lg}) zur Identifikation von Auslassungen und Paraphrasen.
    \item \textbf{Syntaktische Shifts:} Verteilung der Wortarten (Part-of-Speech / POS) und Quantifizierung des Rückgangs von Subordinationen und Passivstrukturen.
\end{itemize}